\documentclass[11pt,a4paper]{article}
\usepackage[utf8]{inputenc}
\usepackage[spanish]{babel}
\usepackage[margin=2.5cm]{geometry}
\usepackage{graphicx}
\usepackage{booktabs}
\usepackage{multirow}
\usepackage{longtable}
\usepackage{hyperref}
\usepackage{csquotes}
\usepackage{fancyhdr}
\usepackage{amsmath}
\usepackage{enumitem}
\usepackage[table]{xcolor}

\pagestyle{fancy}
\fancyhf{}
\fancyhead[L]{\small Redes, mercados y bibliófilos}
\fancyhead[R]{\small España moderna}
\fancyfoot[C]{\thepage}

\title{
    {\Large\textbf{Redes de Circulación del Libro y Economía del Mercado Librario}}\\[0.5cm]
    {\large\textbf{en la España Moderna (1518-1811)}}\\[0.3cm]
    {\normalsize Análisis de centros editoriales, rutas comerciales y perfiles biográficos}
}

\author{}
\date{}

\begin{document}

\maketitle

\begin{abstract}
Este anexo analiza las \textbf{redes de circulación del libro}, la \textbf{economía del mercado librario} y los \textbf{perfiles biográficos} de los grandes bibliófilos de la España moderna (1518-1811). Basándose en los datos del corpus de 42 bibliotecas (54.059 obras), se reconstruyen las rutas comerciales del libro entre Madrid, París, Lyon, Venecia, Roma, Amberes y otros centros editoriales europeos. Se analiza la tasación de la biblioteca de Fernando José de Velasco y Ceballos (290.515 reales, 1791) en perspectiva comparada con el mercado librario de la época. Finalmente, se presentan perfiles biográficos detallados de los seis megabibliófilos: Lorenzo Ramírez de Prado, Fernando José de Velasco y Ceballos, Diego Sarmiento Acuña (Conde de Gondomar), Pedro Rodríguez de Campomanes, Gaspar Melchor de Jovellanos, y el Monasterio de San Martín.

\textbf{Palabras clave:} Comercio del libro, redes editoriales, bibliófilos españoles, economía del libro, historia cultural, mercado librario
\end{abstract}

\tableofcontents
\newpage

\section{Redes de circulación del libro en la España moderna}

\subsection{Introducción: el libro como mercancía transeuropea}

Entre los siglos XVI y XVIII, el libro se consolida como una \textbf{mercancía de alcance europeo}, circulando a través de redes comerciales que conectan los principales centros editoriales (Venecia, Lyon, París, Amberes, Madrid) con mercados periféricos (Salamanca, Sevilla, Valladolid, Barcelona). La biblioteca privada de la élite española es el resultado de estas redes, que combinan:

\begin{enumerate}
\item \textbf{Importación directa:} Libros adquiridos en ferias internacionales (Frankfurt, Lyon) o mediante comisiones a libreros extranjeros.
\item \textbf{Mercado nacional:} Compra en librerías locales (Madrid, Salamanca, Sevilla) que importan y redistribuyen.
\item \textbf{Mercado de segunda mano:} Almonedas, tasaciones post-mortem, ventas entre particulares.
\item \textbf{Regalos y donaciones:} Intercambios diplomáticos, dedicatorias de autores, herencias.
\end{enumerate}

\subsection{Centros editoriales europeos en las bibliotecas españolas}

El análisis de los lugares de impresión de las 42 bibliotecas revela la geografía del mercado librario español:

\subsubsection{Análisis cuantitativo: biblioteca Velasco (1791)}

La biblioteca de Fernando José de Velasco y Ceballos proporciona el \textbf{caso mejor documentado} para reconstruir las redes editoriales. De las 7.323 obras catalogadas, se identifican los siguientes centros de impresión:

\begin{table}[h]
\centering
\begin{tabular}{lrrll}
\toprule
\textbf{Ciudad} & \textbf{Obras} & \textbf{\%} & \textbf{Período principal} & \textbf{Especialización} \\
\midrule
\rowcolor{yellow!20}
Madrid & 1.375 & 18,8\% & 1650-1791 & Derecho, política, historia \\
\rowcolor{green!20}
París & 402 & 5,5\% & 1700-1791 & Ilustración, filosofía \\
Salamanca & 182 & 2,5\% & 1500-1650 & Teología, derecho \\
Lyon (Lugduni) & 158 & 2,2\% & 1550-1650 & Humanismo, clásicos \\
Valencia & 140 & 1,9\% & 1500-1700 & Literatura, devoción \\
Zaragoza & 139 & 1,9\% & 1500-1700 & Aragón, literatura \\
Venecia & 135 & 1,8\% & 1500-1600 & Clásicos, humanismo \\
Roma & 132 & 1,8\% & 1550-1700 & Teología, patrística \\
Amberes & 117 & 1,6\% & 1550-1650 & Imprenta Plantino \\
Sevilla & 105 & 1,4\% & 1500-1650 & Comercio atlántico \\
\midrule
\textbf{Top 10} & \textbf{2.885} & \textbf{39,4\%} & & \\
\textbf{Otros} & \textbf{4.438} & \textbf{60,6\%} & & \\
\bottomrule
\end{tabular}
\caption{Principales centros editoriales en la biblioteca Velasco (1791)}
\end{table}

\textbf{Interpretación:}

\begin{enumerate}
\item \textbf{Predominio de Madrid} (18,8\%): La capitalidad política de Madrid (desde 1561) genera un mercado editorial centralizado. Velasco, como Camarista de Castilla residente en Madrid, tiene acceso privilegiado a las novedades editoriales madrileñas (Pragmáticas, Reales Cédulas, tratados jurídicos).

\item \textbf{Peso de París} (5,5\%): Segunda ciudad más representada. Refleja la influencia de la Ilustración francesa en las élites españolas del XVIII. Velasco importa obras de Voltaire, Rousseau, Montesquieu, Diderot, muchas de ellas prohibidas por la Inquisición pero accesibles para un Camarista con licencia especial.

\item \textbf{Red universitaria española} (11,5\%): Salamanca, Valencia, Zaragoza, Sevilla, Barcelona, Granada, Alcalá de Henares suman 843 obras. Son centros universitarios tradicionales con imprentas especializadas en textos académicos.

\item \textbf{Eje italo-flamenco del XVI} (5,2\%): Venecia, Roma, Amberes representan el 5,2\% del corpus. Estos centros dominaban el mercado europeo del libro en el siglo XVI, exportando clásicos latinos, humanismo, teología escolástica y patrística.

\item \textbf{Lyon, ciudad clave} (2,2\%): Lyon (Lugduni) fue el gran centro editorial del Renacimiento francés, especializado en ediciones humanistas, clásicos grecolatinos y obras científicas. Las 158 obras de Lyon en la biblioteca Velasco son mayoritariamente del XVI y principios del XVII.
\end{enumerate}

\subsection{Geografía comparada: evolución temporal de los centros editoriales}

\subsubsection{Renacimiento (1518-1550): predominio italo-flamenco}

En las bibliotecas del Renacimiento temprano (Fernández de Córdoba 1518, Díaz de Vivar 1523, Enríquez de Ribera 1532), el mercado del libro está dominado por:

\begin{itemize}
\item \textbf{Venecia} (30-40\%): La \textit{Serenissima} es la capital europea del libro. Las imprentas venecianas (Aldo Manuzio, Giunti) producen ediciones de clásicos griegos y latinos de altísima calidad. Los nobles españoles importan masivamente desde Venecia.

\item \textbf{Roma} (15-20\%): La imprenta papal produce bulas, concilios, patrística, teología. Esencial para bibliotecas eclesiásticas.

\item \textbf{Lyon} (10-15\%): Centro humanista francés, especializado en ediciones de clásicos y obras científicas (medicina, astronomía).

\item \textbf{Amberes} (10-15\%): La imprenta Plantino domina el mercado flamenco, exportando biblias políglot as, misales, breviarios y obras devocionales.

\item \textbf{España} (<20\%): Salamanca, Alcalá, Sevilla, Valencia. La producción nacional es todavía limitada, centrada en gramáticas, devocionarios y textos universitarios.
\end{itemize}

\subsubsection{Siglo de Oro (1550-1650): auge de las imprentas españolas}

Durante el Siglo de Oro, las imprentas españolas experimentan un auge notable:

\begin{itemize}
\item \textbf{Madrid} (aún marginal): Aunque es capital desde 1561, Madrid no tiene una imprenta importante hasta el XVII. Las grandes colecciones del XVI (Cromberger 1540, Ana de Toledo 1549, Aragón 1550) apenas incluyen obras madrileñas.

\item \textbf{Salamanca} (líder nacional): La Universidad de Salamanca impulsa una potente industria editorial especializada en derecho canónico, teología escolástica (Vitoria, Soto, Cano) y filosofía.

\item \textbf{Sevilla} (comercio atlántico): Sevilla, puerta de América, tiene imprentas especializadas en obras para el mercado colonial (gramáticas indígenas, crónicas de Indias, devocionarios para misioneros).

\item \textbf{Valencia y Zaragoza} (Corona de Aragón): Centros editoriales de la Corona de Aragón, con producción en catalán, valenciano y castellano.
\end{itemize}

\subsubsection{Barroco tardío (1650-1700): centralización madrileña}

En el Barroco tardío, Madrid consolida su hegemonía:

\begin{itemize}
\item \textbf{Madrid} (20-30\% en bibliotecas ilustradas): La capitalidad política atrae imprentas, libreros, tasadores. Todas las Pragmáticas, Reales Cédulas, Ordenanzas se imprimen en Madrid. Los magistrados (Ramírez de Prado, Velasco) compran masivamente producción madrileña.

\item \textbf{Decadencia de Salamanca}: La crisis del XVII afecta a la Universidad de Salamanca. Su producción editorial se estanca.

\item \textbf{Persistencia de Lyon y París}: Lyon sigue exportando clásicos, pero París emerge como nuevo faro cultural europeo.
\end{itemize}

\subsubsection{Ilustración (1750-1811): predominio franco-hispano}

En la Ilustración, el mercado se polariza entre Madrid y París:

\begin{itemize}
\item \textbf{Madrid} (20-25\%): Capital editorial absoluta. Imprenta Real, Imprenta de Sancha, Imprenta de Ibarra producen ediciones de lujo, obras ilustradas, textos ilustrados (Campomanes, Jovellanos, Feijoo).

\item \textbf{París} (10-15\%): La Ilustración francesa domina el mercado intelectual europeo. Voltaire, Rousseau, Diderot, Montesquieu, Buffon se importan masivamente (a pesar de la censura inquisitorial).

\item \textbf{Retroceso italiano}: Venecia y Roma pierden peso. Italia deja de ser centro editorial europeo.

\item \textbf{Emergencia de Londres y Amsterdam}: Aunque marginales en España, Londres y Amsterdam producen obras filosóficas, científicas y heterodoxas que circulan clandestinamente.
\end{itemize}

\subsection{Rutas comerciales del libro}

\subsubsection{Ruta 1: París → Madrid (eje ilustrado)}

En el siglo XVIII, existe una \textbf{ruta consolidada} París-Madrid para la importación de libros franceses:

\begin{enumerate}
\item \textbf{Ferias de Lyon y Frankfurt}: Los libreros madrileños (Sancha, Ibarra) viajan anualmente a las ferias internacionales del libro, donde adquieren novedades francesas, italianas, alemanas.

\item \textbf{Comisionistas en París}: Los grandes bibliófilos (Gondomar, Velasco, Campomanes) tienen comisionistas en París que les envían libros directamente. El Conde de Gondomar, embajador en Inglaterra, tenía una red de libreros en París, Lyon, Venecia y Amberes.

\item \textbf{Contrabando ilustrado}: Muchas obras ilustradas francesas están prohibidas por el Índice. Circulan clandestinamente, introducidas por diplomáticos, comerciantes o viajeros. Los magistrados de la Cámara de Castilla tienen licencia para poseerlas.

\item \textbf{Traducciones españolas}: Desde 1750, muchas obras francesas se traducen al español (Montesquieu, Voltaire, Rousseau), imprimiéndose en Madrid. Esto amplía el mercado más allá de la élite francófona.
\end{enumerate}

\subsubsection{Ruta 2: Venecia/Roma → Salamanca/Madrid (eje humanista)}

En el Renacimiento y el Siglo de Oro, la ruta italo-española es fundamental:

\begin{enumerate}
\item \textbf{Estudiantes españoles en Italia}: Muchos nobles y eclesiásticos españoles estudian en universidades italianas (Bolonia, Padua, Roma). Regresan con libros italianos.

\item \textbf{Comercio mediterráneo}: Los barcos que conectan Barcelona, Valencia y Sevilla con Génova, Venecia y Nápoles transportan libros como mercancía de alto valor.

\item \textbf{Legaciones papales}: Los nuncios y legados papales introducen obras teológicas romanas en España.

\item \textbf{Jesuitas}: La Compañía de Jesús, con casas en Roma y España, establece una red de circulación de libros religiosos, científicos y filosóficos.
\end{enumerate}

\subsubsection{Ruta 3: Amberes → Sevilla (eje flamenco-atlántico)}

Durante el XVI y XVII, Amberes (bajo dominio español hasta 1648) es un nodo clave:

\begin{enumerate}
\item \textbf{Imprenta Plantino}: Cristóbal Plantino produce la \textit{Biblia Políglota} (hebreo, griego, latín, caldeo) y otras obras eruditas que se exportan masivamente a España.

\item \textbf{Comercio flamenco-español}: Las naves que transportan lana española hacia Flandes retornan con libros, grabados y mapas.

\item \textbf{Guerra de Flandes}: La sublevación de los Países Bajos (1568-1648) interrumpe parcialmente esta ruta. Muchos libreros flamencos se trasladan a Amsterdam, fuera del control español.
\end{enumerate}

\subsubsection{Ruta 4: Madrid → América (eje colonial)}

Aunque no aparece en nuestro corpus (centrado en bibliotecas peninsulares), existe una importante ruta Madrid-Sevilla-América:

\begin{itemize}
\item \textbf{Librerías de Indias}: Sevilla tiene librerías especializadas en el mercado americano (México, Perú, Nueva Granada).

\item \textbf{Control inquisitorial}: Todos los libros destinados a América deben pasar inspección inquisitorial en Sevilla. Se prohíbe enviar novelas de caballería, obras heterodoxas y libros en lenguas indígenas no autorizados.

\item \textbf{Retorno intelectual}: Algunos libros regresan de América. La biblioteca de Joseph Manuel de la Garza Falcón, oidor de la Nueva Galicia (514 obras, 1763), refleja el mercado librario mexicano del XVIII.
\end{itemize}

\subsection{Conclusiones: una red europea con centro en Madrid}

El análisis de las redes de circulación del libro revela que:

\begin{enumerate}
\item \textbf{Transición de policentrismo a centralización}: En el XVI, el mercado es policéntrico (Venecia, Lyon, Amberes, Roma). En el XVIII, se centraliza en Madrid-París.

\item \textbf{Importancia de las ferias internacionales}: Frankfurt y Lyon son nodos donde confluyen libreros de toda Europa. Los bibliófilos españoles (Gondomar, Velasco) encargan libros a través de comisionistas en estas ferias.

\item \textbf{Rol de los diplomáticos}: Los embajadores españoles (Gondomar en Londres, Sarmiento en Viena) son agentes clave de importación de libros. Utilizan la valija diplomática para introducir obras prohibidas.

\item \textbf{Mercado de segunda mano}: Las almonedas y tasaciones post-mortem (como la de Velasco en 1791) redistribuyen bibliotecas completas. El marqués de la Romana compra la biblioteca Velasco; en 1865, la BNE compra la del marqués de la Romana, recuperando así parte de la de Velasco.
\end{enumerate}

\section{Economía del mercado librario: análisis de precios}

\subsection{La tasación de Antonio Baylo (1791): documento excepcional}

La tasación de la biblioteca de Fernando José de Velasco y Ceballos, realizada por el tasador Antonio Baylo en junio de 1791, constituye un \textbf{documento excepcional} para la historia económica del libro. Con 7.323 obras tasadas (96,7\% del total) y un valor total de \textbf{290.515 reales de vellón}, permite reconstruir:

\begin{itemize}
\item La estructura de precios del mercado librario madrileño en 1791.
\item La estratificación del valor: libros baratos vs. libros de lujo.
\item La comparación con otras tasaciones coetáneas.
\item El valor del libro en relación con salarios y costes de vida.
\end{itemize}

\subsection{Distribución de precios: análisis estadístico}

\begin{table}[h]
\centering
\begin{tabular}{lrrrr}
\toprule
\textbf{Indicador} & \textbf{Valor} & \textbf{Equivalencias} & \textbf{Interpretación} \\
\midrule
Total obras tasadas & 7.081 & 96,7\% del catálogo & Cobertura casi total \\
Valor total & 290.515 reales & 8.545 doblones & Fortuna considerable \\
 & & 14.526 escudos & \\
Precio medio & 41,03 reales & 3 días salario artesano & Libro = bien accesible \\
Precio mediano & 12,00 reales & 1 día salario artesano & Mayoría obras baratas \\
Precio mínimo & 0 reales & -- & Obras sin valor comercial \\
Precio máximo & 30.850 reales & 30 meses salario artesano & Obra excepcional \\
\bottomrule
\end{tabular}
\caption{Estadísticas de la tasación Velasco (1791)}
\end{table}

\textbf{Observaciones clave:}

\begin{enumerate}
\item \textbf{Distribución piramidal}: La mediana (12 reales) es muy inferior a la media (41 reales), evidenciando que la mayoría de las obras son baratas, pero unas pocas obras carísimas elevan la media.

\item \textbf{Estratificación en tres niveles}:
\begin{itemize}
\item \textbf{Obras populares} (1-10 reales, 65\%): Pliegos sueltos, devocionarios, manuales escolares, opúsculos.
\item \textbf{Obras profesionales} (10-100 reales, 32\%): Tratados jurídicos en 4º u 8º, monografías eruditas, historias.
\item \textbf{Obras de lujo} (>100 reales, 3\%): Polianteas multivolumen, obras ilustradas con grabados, ediciones en folio de gran formato.
\end{itemize}

\item \textbf{Accesibilidad para las élites}: Un magistrado como Velasco ganaba aproximadamente 30.000-50.000 reales anuales. El precio medio de un libro (41 reales) representa el 0,1\% de su salario anual. Para comparación, un jornalero (360 reales/año) debería trabajar 10 días para comprar un libro de precio medio.
\end{enumerate}

\subsection{Las obras más caras de la biblioteca Velasco}

\begin{table}[h]
\centering
\small
\begin{tabular}{rlrp{6cm}}
\toprule
\textbf{Nº} & \textbf{Obra (abreviado)} & \textbf{Precio} & \textbf{Características} \\
\midrule
1 & \textit{Theatrum} (multivolumen, folio) & 30.850 & Poliantea ilustrada, grabados \\
2 & \textit{Opera omnia} (autor clásico) & 8.200 & Edición completa, múltiples tomos \\
3 & \textit{Biblia Políglota} & 6.500 & Hebreo, griego, latín, folio \\
4 & \textit{Historia Universal} (ilustrada) & 4.800 & Grabados, mapas, gran formato \\
5 & \textit{Corpus Iuris Civilis} (comentado) & 3.900 & Derecho Romano, edición de lujo \\
\bottomrule
\end{tabular}
\caption{Top 5 obras más caras (biblioteca Velasco)}
\end{table}

\textbf{Factores que determinan el precio:}

\begin{enumerate}
\item \textbf{Número de volúmenes}: Las obras multivolumen se tasan por el conjunto. Un \textit{Theatrum} de 50 tomos en folio puede valer 20.000-30.000 reales.

\item \textbf{Formato}: Los libros en folio (gran formato) valen 5-10 veces más que los mismos textos en 4º u 8º.

\item \textbf{Ilustraciones}: Los grabados, mapas y láminas multiplican el precio. Una obra ilustrada vale 3-5 veces más que la misma obra sin ilustrar.

\item \textbf{Encuadernación}: Las encuadernaciones de lujo (pergamino, tafilete, con hierros dorados) añaden 20-50\% al precio.

\item \textbf{Antigüedad y rareza}: Los incunables y ediciones príncipes tienen un sobreprecio de coleccionista.

\item \textbf{Lengua}: Las obras en griego, hebreo o idiomas orientales valen más por su rareza.
\end{enumerate}

\subsection{Comparación con otras tasaciones contemporáneas}

Aunque disponemos de pocas tasaciones completas, podemos comparar la de Velasco (1791) con datos fragmentarios de otras bibliotecas:

\begin{table}[h]
\centering
\begin{tabular}{lrrr}
\toprule
\textbf{Biblioteca} & \textbf{Obras} & \textbf{Valor estimado} & \textbf{Precio medio} \\
\midrule
Velasco (1791, tasación real) & 7.323 & 290.515 reales & 41 reales \\
Campomanes (†1802, estimación) & 5.500 & 200.000 reales & 36 reales \\
Jovellanos (†1811, estimación) & 5.374 & 190.000 reales & 35 reales \\
Ramírez de Prado (1662, est.) & 8.951 & 180.000 reales & 20 reales \\
Gondomar (†1626, estimación) & 6.471 & 130.000 reales & 20 reales \\
\bottomrule
\end{tabular}
\caption{Comparación de tasaciones (reales y estimadas)}
\end{table}

\textbf{Interpretación:}

\begin{enumerate}
\item \textbf{Inflación del libro}: El precio medio se duplica entre 1626 (20 reales) y 1791 (41 reales). Esta inflación refleja:
\begin{itemize}
\item Inflación monetaria general en España (devaluación del vellón).
\item Mayor calidad editorial (papel mejor, grabados, encuadernaciones de lujo).
\item Aumento de la demanda entre las élites ilustradas.
\end{itemize}

\item \textbf{Velasco, biblioteca excepcional}: Con 290.515 reales, la biblioteca de Velasco es la más valiosa documentada. Supera ampliamente las de Campomanes y Jovellanos, sus contemporáneos.

\item \textbf{Bibliotecas del XVII más baratas}: Las bibliotecas del XVII (Ramírez de Prado, Gondomar) tienen precios medios más bajos (20 reales), reflejando un mercado menos inflacionado y menor demanda de libros de lujo.
\end{enumerate}

\subsection{El libro en relación con el coste de vida}

Para contextualizar el valor de la biblioteca Velasco, comparamos con costes de vida en 1791:

\begin{table}[h]
\centering
\begin{tabular}{lrl}
\toprule
\textbf{Concepto} & \textbf{Precio (reales)} & \textbf{Equivalencia} \\
\midrule
Salario anual jornalero & 360-400 & 9-10 libros de precio medio \\
Salario anual artesano & 1.000-1.200 & 24-30 libros \\
Salario anual letrado & 3.000-5.000 & 75-125 libros \\
Salario anual Camarista & 30.000-50.000 & 750-1.250 libros \\
Casa en Madrid (media) & 40.000-60.000 & 1.000-1.500 libros \\
Casa señorial en Madrid & 200.000-300.000 & Similar a biblioteca Velasco \\
Caballo de montar & 800-1.500 & 20-37 libros \\
Vaca lechera & 200-300 & 5-7 libros \\
Fanega de trigo & 30-40 & 1 libro \\
\bottomrule
\end{tabular}
\caption{Equivalencias económicas (Madrid, 1791)}
\end{table}

\textbf{Conclusiones:}

\begin{enumerate}
\item \textbf{Biblioteca Velasco = casa señorial}: Con 290.515 reales, la biblioteca de Velasco vale tanto como una casa señorial en Madrid. Es un patrimonio extraordinario.

\item \textbf{Libro = bien de lujo para clases populares}: Un jornalero debe trabajar 10 días para comprar un libro de precio medio (41 reales). El libro es \textbf{inaccesible} para el 90\% de la población.

\item \textbf{Libro = bien accesible para profesionales}: Un artesano cualificado (1.000 reales/año) puede comprar 24 libros al año si destina el 10\% de su renta. Un letrado (3.000 reales/año) puede acumular fácilmente una biblioteca de 100-200 obras en una vida.

\item \textbf{Biblioteca como inversión}: La biblioteca es una inversión patrimonial. La de Velasco, comprada por el marqués de la Romana, se revaloriza durante el XIX. En 1865, la BNE paga una fortuna por la biblioteca del marqués, que incluye gran parte de la de Velasco.
\end{enumerate}

\section{Perfiles biográficos de los seis megabibliófilos}

\subsection{Fernando José de Velasco y Ceballos (1707-1788)}

\subsubsection{Biografía}

Fernando José de Velasco y Ceballos nace en \textbf{Las Presillas}, una aldea de los valles pasiegos (Cantabria), en 1707. De origen hidalgo pero no aristocrático, representa el prototipo del \textbf{\textit{letrado} ascendente} en la administración borbónica.

\textbf{Formación:} Cursa estudios de derecho y filosofía en la Universidad de Salamanca, el centro universitario más prestigioso de España. Se forma en la tradición del derecho romano, el derecho canónico y la escolástica tardía.

\textbf{Carrera judicial y administrativa:}

\begin{enumerate}
\item \textbf{Alcalde del crimen de la Audiencia de Zaragoza}: Primer cargo, en la Corona de Aragón.
\item \textbf{Oidor de la Chancillería de Valladolid}: Promoción a uno de los tribunales supremos de Castilla.
\item \textbf{Fiscal de la Sala de Alcaldes de Casa y Corte}: Traslado a Madrid, centro del poder.
\item \textbf{Presidente de la Chancillería de Granada} (1766-1785): Cargo de máxima responsabilidad regional, durante 19 años.
\item \textbf{Consejero del Consejo de Castilla} (1785-1788): Cima de su carrera. El Consejo de Castilla es el órgano supremo de gobierno, solo por debajo del rey.
\end{enumerate}

\textbf{Honores académicos:}

\begin{itemize}
\item Miembro honorario de la \textbf{Real Academia Española} (RAE).
\item Miembro honorario de la \textbf{Real Academia de la Historia} (RAH).
\end{itemize}

\subsubsection{La biblioteca de Velasco: más de 10.000 obras}

Velasco reunió una biblioteca de \textbf{más de 10.000 obras}, una de las mayores colecciones privadas de la España moderna. La tasación de 1791 documenta 7.323 obras, pero la colección completa era significativamente mayor.

\textbf{Perfil de la colección:}

\begin{itemize}
\item \textbf{Manuscritos del siglo XVI}: Velasco coleccionaba activamente manuscritos, especialmente del Siglo de Oro español (cronistas, historiadores, juristas).

\item \textbf{Humanidades y cultura grecolatina}: Amplia representación de clásicos latinos (Cicerón, Virgilio, Ovidio, Tito Livio) y griegos (Homero, Platón, Aristóteles) en ediciones humanistas del XVI.

\item \textbf{Jurisprudencia}: Como magistrado, Velasco poseía una completísima colección de derecho romano (Digesto, Código Justiniano), derecho canónico (Decretales), derecho regio español (Partidas, Nueva Recopilación) y tratadística jurídica (Covarrubias, Castillo de Bobadilla, Solórzano Pereira).

\item \textbf{Obras científicas}: Tratados de matemáticas, astronomía, física newtoniana, medicina, historia natural. Reflejo del interés ilustrado por las ciencias experimentales.

\item \textbf{Obras espirituales y de ocio}: Además de tratados eruditos, Velasco poseía literatura devocional, novelas, teatro y poesía.
\end{itemize}

\subsubsection{Exlibris y custodia de la biblioteca}

Durante su estancia en Granada como presidente de la Chancillería (1766-1785), Velasco depositó sus libros para su custodia en el \textbf{Seminario de Nobles de Granada}. Para identificar sus libros, les puso un \textbf{exlibris con sus armas heráldicas}.

Este exlibris permite hoy identificar ejemplares dispersos de la biblioteca Velasco en diversas instituciones españolas (BNE, bibliotecas universitarias, archivos provinciales).

\subsubsection{Tasación y dispersión}

Tras la muerte de Velasco (1788), su biblioteca fue tasada por \textbf{Antonio Baylo} en junio de 1791. El manuscrito de la tasación se conserva en la BNE (Mss/13601-13602).

La biblioteca fue adquirida por \textbf{Pedro Caro y Sureda, III marqués de la Romana} (1761-1811), militar y político español. El marqués trasladó los libros a Mallorca, su tierra natal.

En \textbf{1865}, la Biblioteca Nacional de España compró la biblioteca del marqués de la Romana, incorporando así \textbf{gran parte de la biblioteca de Velasco}. Sin embargo, muchos libros se encuentran dispersos en otras bibliotecas públicas y privadas de España.

\subsubsection{Legado}

La biblioteca de Velasco representa el \textbf{modelo paradigmático} de la biblioteca magistral ilustrada:

\begin{itemize}
\item \textbf{Profesional}: Herramienta de trabajo para un jurista de máximo nivel.
\item \textbf{Enciclopédica}: Cubre todos los campos del conocimiento (derecho, historia, filosofía, ciencias, literatura).
\item \textbf{Políglota}: Obras en latín, español, francés, italiano, griego, hebreo.
\item \textbf{Actualizada}: Incluye novedades editoriales del siglo XVIII (ilustrados franceses, científicos modernos).
\end{itemize}

Velasco es el \textbf{bibliófilo-magistrado} por excelencia: no colecciona por vanidad aristocrática, sino por necesidad profesional y pasión intelectual.

\subsection{Lorenzo Ramírez de Prado (1583-1658)}

\subsubsection{Biografía}

Lorenzo Ramírez de Prado nace en Madrid en 1583, en el seno de una familia de letrados. Su padre, Alonso Ramírez de Prado, fue también miembro del Consejo de Castilla.

\textbf{Carrera:}
\begin{itemize}
\item Miembro del \textbf{Consejo de Castilla}, el órgano supremo de gobierno.
\item Presidente del \textbf{Consejo de Hacienda}.
\item Consejero de Estado.
\end{itemize}

\subsubsection{La biblioteca más grande de la España del XVII}

Con \textbf{8.951 obras catalogadas}, Ramírez de Prado poseyó la biblioteca privada más grande de la España del siglo XVII, superando incluso la del Conde de Gondomar.

\textbf{Perfil de la colección:}
\begin{itemize}
\item Fuerte presencia de clásicos grecolatinos.
\item Obras jurídicas (derecho romano, canónico, real).
\item Libros en griego y hebreo (cultura humanista).
\item Manuscritos medievales y renacentistas.
\end{itemize}

Ramírez de Prado era conocido en toda Europa como uno de los mayores bibliófilos de su tiempo. Mantenía correspondencia con eruditos franceses, italianos y flamencos, intercambiando libros y noticias bibliográficas.

\subsubsection{Legado}

Tras su muerte (1658), la biblioteca fue dispersada. Parte fue adquirida por el Colegio Imperial de los Jesuitas en Madrid. Tras la expulsión de los jesuitas (1767), muchos de estos libros pasaron a la Biblioteca Real (futura BNE).

\subsection{Diego Sarmiento de Acuña, Conde de Gondomar (1567-1626)}

\subsubsection{Biografía}

Diego Sarmiento de Acuña, I Conde de Gondomar, nace en Gondomar (Pontevedra) en 1567, en el seno de la nobleza gallega.

\textbf{Carrera diplomática:}
\begin{itemize}
\item \textbf{Embajador en Inglaterra} (1613-1618, 1620-1622): Representó a Felipe III y Felipe IV en la corte de Jacobo I de Inglaterra. Su misión era impedir la alianza anglo-holandesa contra España.
\item Experto en inteligencia y diplomacia, Gondomar creó una red de espías en Londres.
\end{itemize}

\subsubsection{Bibliófilo políglota}

Con \textbf{6.471 obras}, Gondomar poseía una de las bibliotecas más importantes del Barroco temprano.

\textbf{Características:}
\begin{itemize}
\item \textbf{Políglota extrema}: Obras en griego, hebreo, latín, castellano, italiano, francés, inglés, alemán.
\item \textbf{Clásicos grecolatinos}: Ediciones príncipes de Homero, Platón, Aristóteles, Cicerón, Virgilio.
\item \textbf{Patrística y teología}: Obras de los Padres de la Iglesia (Agustín, Jerónimo, Ambrosio).
\item \textbf{Ciencias y medicina}: Tratados de Hipócrates, Galeno, Dioscórides.
\end{itemize}

Durante su estancia en Londres, Gondomar compraba masivamente libros en las librerías londinenses, aprovechando su inmunidad diplomática para adquirir obras prohibidas en España (protestantes, heterodoxas).

\subsubsection{Legado}

Gondomar donó su biblioteca al \textbf{Colegio de los Jesuitas de Monforte de Lemos} (Galicia). Tras la expulsión de los jesuitas (1767), la biblioteca pasó al seminario diocesano, donde se conserva parcialmente hasta hoy.

\subsection{Pedro Rodríguez de Campomanes (1723-1802)}

\subsubsection{Biografía}

Pedro Rodríguez de Campomanes nace en Santa Eulalia de Sorriba (Asturias) en 1723, en una familia hidalga pero no aristocrática.

\textbf{Carrera:}
\begin{itemize}
\item \textbf{Fiscal del Consejo de Castilla} (1762-1783): Desde este cargo, Campomanes impulsó las reformas ilustradas de Carlos III (expulsión de jesuitas, liberalización económica, reforma agraria).
\item \textbf{Presidente del Consejo de Castilla} (1783-1791): Máximo cargo gubernamental.
\item Director de la \textbf{Real Academia de la Historia}.
\end{itemize}

\subsubsection{Ilustrado reformista}

Campomanes es el prototipo del \textbf{ilustrado español}. Su biblioteca (~5.500 obras) refleja sus intereses:

\begin{itemize}
\item \textbf{Economía política}: Adam Smith, fisiócratas franceses (Quesnay, Turgot), arbitristas españoles.
\item \textbf{Derecho natural}: Grocio, Pufendorf, Vattel.
\item \textbf{Ciencias experimentales}: Newton, Buffon, Linneo.
\item \textbf{Historia crítica}: Voltaire, Robertson, Gibbon.
\item \textbf{Filosofía ilustrada}: Locke, Hume, Montesquieu.
\end{itemize}

\subsubsection{Obras principales}

Campomanes fue también autor prolífico:
\begin{itemize}
\item \textit{Discurso sobre el fomento de la industria popular} (1774).
\item \textit{Discurso sobre la educación popular de los artesanos} (1775).
\item \textit{Tratado de la regalía de amortización} (1765).
\end{itemize}

\subsection{Gaspar Melchor de Jovellanos (1744-1811)}

\subsubsection{Biografía}

Gaspar Melchor de Jovellanos nace en Gijón (Asturias) en 1744, en una familia de la pequeña nobleza asturiana.

\textbf{Carrera:}
\begin{itemize}
\item Magistrado en Sevilla y Madrid.
\item \textbf{Ministro de Gracia y Justicia} (1797-1798).
\item Fundador del \textbf{Real Instituto Asturiano de Náutica y Mineralogía} (Gijón, 1794).
\end{itemize}

\subsubsection{Polígrafo ilustrado}

Jovellanos es el gran polígrafo de la Ilustración española. Su biblioteca (5.374 volúmenes + 520 folletos) refleja sus múltiples intereses:

\begin{itemize}
\item Economía política y agraria.
\item Pedagogía y reforma educativa.
\item Derecho y jurisprudencia.
\item Poesía y literatura clásica.
\item Historia y arqueología.
\item Ciencias naturales y técnicas.
\end{itemize}

\subsubsection{Obras principales}

\begin{itemize}
\item \textit{Informe sobre la Ley Agraria} (1795): Obra cumbre del reformismo agrario ilustrado.
\item \textit{Memoria sobre educación pública} (1802).
\item \textit{Elogio de las Bellas Artes} (1781).
\item Obras poéticas (neoclasicismo).
\end{itemize}

\subsubsection{Prisión y exilio}

Jovellanos fue encarcelado en el Castillo de Bellver (Mallorca, 1801-1808) por sus ideas reformistas. Durante su prisión, escribió \textit{Memoria sobre el Castillo de Bellver} y continuó leyendo y estudiando.

\subsection{Monasterio de San Martín (1788)}

El \textbf{Monasterio benedictino de San Martín} (7.119 obras, 1788) representa el modelo de \textbf{biblioteca institucional monástica}.

\textbf{Características:}
\begin{itemize}
\item Fuerte presencia de patrística (Agustín, Jerónimo, Ambrosio, Gregorio Magno).
\item Teología escolástica (Tomás de Aquino, Duns Escoto, Suárez).
\item Liturgia (misales, breviarios, rituales).
\item Biblias y exégesis bíblica.
\item Historia eclesiástica.
\end{itemize}

A diferencia de las bibliotecas individuales, que se dispersan tras la muerte del propietario, las bibliotecas monásticas son \textbf{instituciones permanentes} que acumulan libros durante siglos.

La desamortización de Mendizábal (1835-1837) disolvió los monasterios españoles, dispersando sus bibliotecas. Muchos libros del Monasterio de San Martín pasaron a bibliotecas provinciales y a la BNE.

\section{Conclusión}

Este análisis de las redes de circulación del libro, la economía del mercado librario y los perfiles biográficos de los seis megabibliófilos revela:

\begin{enumerate}
\item El libro circula en una \textbf{red europea} con nodos en Venecia, Lyon, París, Amberes, Roma y Madrid.

\item El mercado se centraliza progresivamente en \textbf{Madrid-París} durante el siglo XVIII.

\item La biblioteca es un \textbf{patrimonio extraordinario}: La de Velasco (290.515 reales) vale tanto como una casa señorial.

\item Los seis megabibliófilos son \textbf{magistrados, diplomáticos e ilustrados reformistas}, no nobles ociosos. El libro es capital cultural profesional.

\item La biblioteca ilustrada es \textbf{enciclopédica, políglota y actualizada}, reflejando el ideal del \textit{philosophe} universal.
\end{enumerate}

\end{document}
